\documentclass[11pt,a4paper]{article}
\usepackage[T1]{fontenc}
\usepackage[utf8]{inputenc}
\usepackage[spanish]{babel}
\usepackage{amsmath,amssymb,amsthm,mathtools}
\usepackage[margin=25mm]{geometry}
\usepackage[hidelinks]{hyperref}
\newtheorem{proposition}{Proposición}
\newtheorem{theorem}{Teorema}
\newcommand{\F}{\mathbb F}
\newcommand{\Z}{\mathbb Z}
\newcommand{\Q}{\mathbb Q}
\title{Composición generativa, incidencia excepcional\protect\\
y recuperación de la constante holográfica\protect\\
de acoplamiento aritmético--geométrico}
\author{Nota técnica de continuidad y formalización}
\date{18 de septiembre de 2026}
\begin{document}
\maketitle

\begin{abstract}
Se reúnen los dominios y las composiciones que enlazan las semillas
APP--TRIT--TPK con la incidencia excepcional y con la recuperación del registro
dodecafásico. La exposición conserva la distinción entre selección de una
historia, publicación de una incidencia, recuperación de un registro y lectura
de la constante asociada. Incluye la selección de la marca reticular, la
procedencia del coeficiente radial $4$ y de su norma $54$, un recuperador
explícito mediante incidencias de Steiner y los puntos de integración con
Lean. No sustituye el tratado integral ni declara ejecutado un productor
terminal que los módulos de esta nota no contienen.
\end{abstract}

\section{Tesis rectora y genealogía de la construcción}

El núcleo operativo de Holografía Modular Triádica y Mecánica Dimensional
comprende conjuntamente la APP, el TRIT, las reglas dinámicas del TPK, el
estado enriquecido y la estructura discreta del continuo que resulta de su
extensión. La expresión «núcleo mínimo» designa esta composición íntegra;
no designa una APP aislada de su dinámica ni un catálogo de valores
numéricos. Su orden causal se mantiene en toda la nota:
\[
\begin{aligned}
 \mathrm{APP}&\longrightarrow\mathrm{TRIT}\longrightarrow\mathrm{TPK}\\
 &\longrightarrow\text{estado enriquecido}\\
 &\longrightarrow\text{estructura discreta conjunta del continuo}.
\end{aligned}
\]
Las hojas de suma y producto actúan sobre las mismas marcas APP. Para
$p,q\in\{1,\ldots,9\}$, sus lecturas conservan las descomposiciones
\[
 p+q=r_\Sigma+9q_\Sigma,\qquad
 pq=r_\Pi+9q_\Pi,\qquad 1\leq r_\Sigma,r_\Pi\leq9.
\]
La reducción no elimina el cociente. La aritmética módulo nueve conserva
el nueve como marca de frontera en la representación anterior; la lectura
por tríadas utiliza la base mil y conserva su acarreo. Estos dos niveles
se componen mediante las operaciones del TPK, sin identificar sus cocientes
ni borrar el registro que los relaciona.

El TRIT determina el régimen de lectura, la orientación topológica y la
relación temporal de la operación. El TPK selecciona, transporta,
pliega y despliega, actualiza el registro y prolonga las rutas compatibles.
El estado conserva régimen trítico, orientación, hoja, ruta, acarreo,
frontera y memoria. La generación del catálogo y la selección regional se
efectúan sobre esa información.

En este desarrollo, $\pi$, $\varphi$ y $e$ designan también regiones
generadas de cierre, autoescala y propagación. Sus palabras, sus soportes y
sus relaciones operatorias forman parte de la estructura antes de la
evaluación numérica. Las correspondientes coordenadas numéricas son
lecturas de esas regiones; sus valores convencionales no seleccionan las
semillas por comparación decimal. La utilización posterior de una
coordenada ya generada conserva esta genealogía.

La doble proyección actúa sobre el mismo estado enriquecido. La lectura
arquimediana contrae cilindros compatibles y produce el valor real; la
lectura excepcional conserva la incidencia de frontera a través de las
realizaciones Paley--Witt--Golay--Leech y de su prolongación hacia
$V^\natural$, Monster y Moonshine. Las realizaciones dimensionales y la
dualidad T del desarrollo hacia teoría M pertenecen a la arquitectura del
corpus, con sus mapas y dominios propios. La recta real es la publicación
de esta construcción, no un conjunto de valores objetivo que se introduzca
para seleccionar retrospectivamente las rutas. Los módulos Lean adjuntos
formalizan el tramo regional--incidencial--reticular y la recuperación
finita descritos más abajo; su compilación no se presenta como una nueva
formalización de todas esas realizaciones posteriores.

El registro K y su lectura periódica $\kappa_{\rm ag}$ se sitúan dentro de
esta misma cadena. Su centralidad consiste en enlazar las lecturas
aritméticas, la incidencia y la reconstrucción del estado leído. La
singularidad de una selección se enuncia siempre con las condiciones
conjuntas que la producen: la invertibilidad de un lector aislado no
sustituye la selección del registro. La composición con las regiones
$\pi$, $\varphi$, $e$ determina la lectura de $\alpha$ mediante la ecuación
de acarreo que se expone en la sección correspondiente.

La cadena de publicaciones finitas y su prolongación es
\[
 w_6\longrightarrow w_{12}\longrightarrow w_{18}
 \longrightarrow w_{24}\longrightarrow w_{30}
 \longrightarrow R_{36}\longrightarrow G_9.
\]
Un retorno de fase no identifica los estados enriquecidos anterior y
posterior. La unidad U008 lo exhibe con bloques visibles iguales que conservan
preestados y cocientes distintos. La memoria que distingue esas realizaciones
no puede omitirse al evaluar las ventanas dodecafásicas.

\section{Operadores de elevación y selección regional}

Sobre palabras fila de $\F_3^6$, los registros de transición proporcionan
dos sistemas
\[
 X_0L_0=Y_0,\qquad X_1L_1=Y_1.
\]
La invertibilidad de $X_0$ y $X_1$ determina unívocamente los dos operadores.
Esta unicidad tiene por incógnita la matriz de elevación; las transiciones
que forman $X_j,Y_j$ conservan su procedencia del estado. El calendario común
seleccionado es $(0,0,1,1)$. Las semillas orientadas producidas son
\[
 w_6^\pi=010211,\qquad w_6^e=201101,\qquad
 w_6^\varphi=121200.
\]
La recurrencia $b_{k+1}^\chi=b_k^\chi L_{\epsilon_k}$ produce
\begin{align*}
 w_{30}^{\pi}&=010211|012222|010211|002111|110221,\\
 w_{30}^{e}&=201101|121221|102011|012222|102011,\\
 w_{30}^{\varphi}&=121200|112202|121020|010210|010200.
\end{align*}
Se multiplica cada bloque de seis trits, no la palabra concatenada.

La frontera $R_{36}$ conserva el eje, el agregado, la saturación, la ocupación
y la neutralidad dual. La orientación distingue la terna
\[
 R_{36}=\begin{pmatrix}222220\\021222\\102011\end{pmatrix}
\]
de su realización especular. Este resultado no se reemplaza por la repetición
periódica de una palabra de treinta posiciones.

\section{Incidencia de Paley--Witt y selección de la marca}

En la carta del artículo, el codificador es
\[
 \gamma(w)=(w,wA_W),\qquad
 A_W=\begin{pmatrix}
 0&1&1&1&1&1\\1&0&1&2&2&1\\1&1&0&1&2&2\\
 1&2&1&0&1&2\\1&2&2&1&0&1\\1&1&2&2&1&0
 \end{pmatrix}.
\]
La imagen es el código ternario de dimensión seis y longitud doce del
desarrollo excepcional. Sus 264 palabras de peso seis proporcionan 132
hexadas al identificar soportes iguales.

La aplicación del codificador a las semillas seleccionadas produce
\begin{align*}
 P_\pi&=\{2,4,5,6,7,8,9,10,11\},\\
 H_e&=\{1,3,4,6,9,10\},\\
 H_\varphi&=\{1,2,3,4,7,9\}.
\end{align*}
La palabra APP de la carta, $022020$, tiene soporte
$H_{\rm APP}=\{2,3,5,10,11,12\}$. Su pertenencia al catálogo producido está
formalizada; esa pertenencia no se identifica con un teorema adicional de
selección única de la palabra.

La cuaterna se produce sin recibir los doce bloques de K:
\[
 B=H_e\cap P_\pi=\{4,6,9,10\}.
\]
En la imagen del codificador, la regla
\[
 B\subset H\subset P_\pi,\qquad
 (|H\cap H_e|,|H\cap H_\varphi|,|H\cap H_{\rm APP}|)=(4,3,2)
\]
selecciona una única hexada
\[
 H^-_\alpha=\{4,5,6,7,9,10\}.
\]
El origen resulta de otra operación incidencial:
\[
 (H_e\cap H_\varphi)\setminus(H^-_\alpha\cup H_{\rm APP})=\{1\}.
\]
La elevación hexada--octada se realiza posteriormente en la dodecada y la
carta binaria declaradas. Su unicidad concierne a la octada que eleva una
hexada; no cambia el dominio de la aplicación en una selección de doce
enteros.

\section{Coeficiente radial, norma y vecino reticular}

El código conserva el pegado de doce componentes $A_2$. En la cámara radial
positiva del artículo se escribe
\[
 v(a)=\sum_{i=1}^{12}a_i\rho_i,
 \qquad a_i=1+3b_i,\quad b_i\in\Z_{\geq0},
 \qquad \rho_i^2=2.
\]
La condición del vecino de índice tres es $v(a)^2\equiv0\pmod{18}$.
Como
\[
 v(a)^2=2\sum_i(1+3b_i)^2,
\]
equivale a $\sum_i b_i\equiv1\pmod3$. La norma mínima en esta cámara
se obtiene con un solo $b_i=1$. La marca anterior fija su posición. Así,
\[
 v_\alpha=4\rho_1+\rho_2+\cdots+\rho_{12},
 \qquad v_\alpha^2=2(4^2+11)=54.
\]
El coeficiente $4=1+3$ y la norma $54$ tienen aquí funciones tipadas:
elevación radial mínima y norma del vector de vecino. La minimalidad se
refiere a la cámara positiva declarada. No se confunde esta norma con una
duración cronológica de 54 pasos.

Para el pegado $N=N(C_W)$ se define
\[
 N_{\alpha,0}=\{x\in N:\langle x,v_\alpha\rangle\equiv0\pmod3\},
 \qquad N_\alpha=N_{\alpha,0}+\Z\frac{v_\alpha}{3}.
\]
La composición formalizada conserva positividad, paridad, integralidad,
autodualidad, rango 24 y mínimo 4. Estas propiedades se obtienen de la marca y
del código, sin introducir K en el módulo reunido.

\section{Evaluación del registro y reconstrucción de Hadamard}

El registro de visitas y trazas conserva los recuentos $A_m,C_m,V_m$ por
ventana. En la presentación incidencial del artículo,
\[
 z_m=100[A_m]_{10}+10[C_m]_{10}+[V_m]_{10},
\]
conservando separadamente los cocientes y la procedencia. Sobre registros
$z\in\Z^{12}$, con índices módulo doce, el lector transversal es
\[
 \mathcal E(z)=(D_3z,D_4z,Q(z)),\qquad
 (D_jz)_m=z_m-z_{m+j},\qquad Q(z)=\sum_mz_m.
\]
Los pasos tres y cuatro corresponden a la estructura cíclica de los lectores
90/120. No se insertan como canales numéricos independientes al definir una
ruta del productor.

\begin{proposition}[Unicidad de la reconstrucción transversal]
Si $\mathcal E(z)=\mathcal E(z')$, entonces $z=z'$.
\end{proposition}
\begin{proof}
La diferencia $d=z-z'$ satisface $d_m=d_{m+3}=d_{m+4}$. Los pasos tres y
cuatro generan el ciclo de doce posiciones, por lo que $d$ es constante.
La igualdad de cargas da $12d_1=0$ y anula esa constante.
\end{proof}

En la realización armónica intercalada, $H_{12}^2=4\operatorname{Id}$.
Un registro firmado U y su publicación K satisfacen
\[
 4K=H_{12}U.
\]
El registro U determina un único K en el espacio racional mediante
$K=H_{12}U/4$. Su pertenencia a la red entera se comprueba en el dominio
integral del lector. La selección de U por el estado terminal y la inversión
de esta ecuación se conservan como
dos operaciones de la misma composición, no como demostraciones
intercambiables.

\section{Recuperación íntegra por incidencias de Steiner}

Sea $\mathcal H$ la familia de 132 hexadas y sea
$I:\Q^{12}\longrightarrow\Q^{132}$ el lector
\[
 (Iv)(H)=\sum_{i\in H}v_i.
\]
Cada posición pertenece a 66 hexadas y cada pareja distinta a 30; por tanto
\[
 I^{\mathsf T}I=36\operatorname{Id}+30J.
\]
\begin{theorem}[Recuperador del registro completo]
Para $y=Iv$ se tiene, en cada posición,
\[
 \boxed{396v_i=11\sum_{H\ni i}y_H-5\sum_{H\in\mathcal H}y_H.}
\]
En consecuencia, dos registros con las mismas incidencias ponderadas en
todas las hexadas tienen los mismos doce componentes.
\end{theorem}
\begin{proof}
Sea $S=\sum_iv_i$. Contando las apariciones de cada componente,
\[
 \sum_{H\ni i}y_H=36v_i+30S,\qquad
 \sum_Hy_H=66S.
\]
La combinación $11(36v_i+30S)-5(66S)$ es $396v_i$.
\end{proof}

El certificado Lean adjunto no necesita recibir esta identidad de Gram
como una hipótesis. Selecciona doce hexadas testigo de la imagen del
codificador, comprueba su matriz de incidencia
$B:\Q^{12}\longrightarrow\Q^{12}$ y una matriz racional C,
y demuestra por evaluación exacta
\[
 CB=BC=\operatorname{Id},\qquad C=M/18
\]
para una matriz entera M explícita en el código. Se obtienen las identidades
$\operatorname{recover}(\operatorname{read}(v))=v$ y
$\operatorname{read}(\operatorname{recover}(y))=y$ para esa carta de doce
lecturas. Para el lector completo de 132 coordenadas, la recuperación es
inversa sobre su imagen $I[\Q^{12}]$. La selección de esa base
no consulta K ni sus canales. Las lecturas ponderadas que recibe el inversor
conservan, sin embargo, su procedencia; el inversor no fabrica sus valores.

La lectura racional periódica de la constante se obtiene, una vez producido
el registro ordenado con $K_m\in\{0,\ldots,999\}$ y doce posiciones, mediante
\[
 \kappa_{\rm ag}=\frac{\sum_{m=1}^{12}K_m1000^{12-m}}{1000^{12}-1}.
\]
Esta publicación conserva el orden y permite recuperar los doce bloques en
la base y longitud fijadas, incluidos sus ceros iniciales.

\section{Acoplamiento regional y lectura de la estructura fina}

Sean $P_j,E_j,\Phi_j$ las tríadas publicadas por las regiones de cierre,
propagación y autoescala del mismo desarrollo. Se prolonga el registro
ordenado con su periodo dodecafásico, $K_{j+12}=K_j$. La ecuación posicional
del artículo es
\[
 Z_j=P_j+E_j-\Phi_j-K_j,\qquad s_j=Z_j+c_{j+1},
\]
\[
 A_j=s_j-1000\left\lfloor\frac{s_j}{1000}\right\rfloor,
 \qquad c_j=\left\lfloor\frac{s_j}{1000}\right\rfloor.
\]
Aquí $A_j$ es el bloque de salida de la lectura de la estructura fina,
no el ángulo físico posterior que en otras secciones se denota también
por A. La condición de frontera del acarreo se conserva al pasar de una
profundidad a otra. En particular, fijar arbitrariamente un acarreo terminal
nulo en una ventana finita no reemplaza la compatibilidad de la lectura
infinita.

La ecuación local reúne efectivamente las cuatro lecturas y su registro:
\[
 \boxed{P_j+E_j-\Phi_j-K_j-A_j=1000c_j-c_{j+1}.}
\]
Si $\iota_N(B)=\sum_{j=1}^{N} B_j1000^{N-j}$, su suma telescópica da
\[
\begin{aligned}
 &\iota_N(P)+\iota_N(E)-\iota_N(\Phi)-\iota_N(K)-\iota_N(A)\\
 &\hspace{30mm}=1000^Nc_1-c_{N+1}.
\end{aligned}
\]
Esta identidad precisa en qué sentido el registro complementa las
coordenadas regionales en la publicación de $\alpha$: transporta una
contribución ordenada que debe componerse con la diferencia regional y con
ambos extremos del acarreo. La segunda vía analítica de la fuente conserva
sus propias recurrencias y su identificación mediante cilindros compatibles.
La identidad finita aquí reproducida no sustituye aquella demostración.

El origen de los coeficientes queda separado: mil es la base de la tríada,
doce es el periodo del registro dodecafásico y los signos corresponden a la
composición regional declarada. Ningún valor metrológico de $\alpha$ actúa
como entrada de estas ecuaciones.

\section{Definiciones de los antecedentes del selector terminal}

Los rótulos de archivo W24, S8 y N69 se conservan únicamente como alias
de procedencia. Los objetos matemáticos correspondientes se distinguen por
su construcción y por su función en la composición.

\subsection{Calendario de firmas regionales y reloj de conversión}

El calendario reúne, en cada tiempo de bloque, las tres bandas ternarias
regionales y sus firmas: suma residual, diferencia orientada, cociente,
ocupación, carga visible y carga dual. El reloj de publicación decimal es
\[
 d_t=\max\{k\in\mathbb N:1000^k\leq3^{6(t+1)}\}.
\]
Tiempo de bloque y profundidad decimal son coordenadas distintas. La
operación que calcula las firmas actúa columna a columna sobre las bandas
generadas. Las cien filas del archivo N69 constituyen una ventana de esta
publicación, no su definición ni una fuente independiente de sus valores.
La regeneración regional y la igualdad con esas filas están formalizadas
en \path{GeneratedN69Rows.lean}.

\subsection{Descriptor inicial por retorno nonádico y lectura fase--carga}

Denotamos por $\Omega_{24}$ el descriptor inicial de veinticuatro trits
del registro. Se distingue de $w_{24}^{\chi}$, que es una etapa de la
biografía de una región. La igualdad de longitudes no identifica esos
objetos.

Sea $t_*$ el primer tiempo en que $d_{t+1}=d_t$ y sea $r=t_*-9$.
En la ventana considerada, $t_*=20$ y $r=11$. Con suma componente a
componente en $\F_3$, se conservan los defectos opuestos de semibanda
\[
 \delta_+=000111,\qquad\delta_-=000222,\qquad
 \delta_++\delta_-=0.
\]
El descriptor concatena cuatro contribuciones tipadas:
\begin{enumerate}
\item La banda de propagación en $t_*+1$, corregida por $\delta_-$.
\item Las dos colas de tres trits de $\delta_-$ y $\delta_+$, en ese orden.
\item La suma de las bandas de cierre y propagación en $r+1$.
\item La lectura fase--carga de las bandas en $r$, con fase
$(r-1)\bmod3$ y carga dual en la carta regional declarada.
\end{enumerate}
Los coeficientes del último lector son $+1$ cuando la carga del canal
coincide con la fase y $-1$ en caso contrario. La construcción tiene
longitud $6+3+3+6+6=24$ y produce
\[
 \Omega_{24}=020022222111211201021101.
\]
La cadena de cifras es el resultado del descriptor, no su definición.
Los defectos orientados, la carta de carga y la regla de lectura conservan
su procedencia en el propietario del carácter fase--dual. El primer avance
nulo selecciona el tiempo sin introducir el literal veinte como selector.
El módulo \path{RegionalW24.lean} está compilado y su teorema
\path{regionalW24_eq_input} identifica esta producción con el descriptor
utilizado por el selector. Sus cinco consultas de axiomas sólo declaran
\path{propext}, \path{Classical.choice} y \path{Quot.sound}; no utiliza
evaluación nativa para esta reconstrucción.

\subsection{Repertorio terminal de bloques ternarios}

El repertorio $\mathcal S_{\rm term}\subset\F_3^6$ conserva ocho bloques
distintos, sin fijar su orden. El selector examina sus ordenaciones y sus
cilindros. Su alias histórico es S8. La fuente de procedencia extrae este
repertorio de la segmentación terminal de la palabra de setenta y dos
trits y coteja sus realizaciones mediante los lectores comparativo y afín.
La generación de una familia de bloques, la pertenencia de cada bloque y
la selección exacta del repertorio conservan enunciados distintos. El
recibo del selector no se utiliza para borrar esta dependencia anterior.

\section{Selector heterotípico terminal y continuidad de la implementación}

El desarrollo histórico de cierre afín--monodrómico contiene un selector más
fuerte que la sola cuaterna. Desde el descriptor inicial, el repertorio
terminal no ordenado y el panel de fronteras, construye 19.446 candidatos.
La cara excepcional y el
cambio de tipo de lector en la órbita $4\to8\to12$ seleccionan el registro
\[
 (234,543,140,729,659,824,621,58,914,794,146,601).
\]
Ese resultado conserva el descriptor y el repertorio como antecedentes. Su procedencia no se
sustituye por los valores finales del registro, por los canales esperados
ni por una inversión retrospectiva. La presente nota conserva ese enlace
como propietario del selector; sus dos módulos Lean certifican la composición
excepcional y el recuperador, no toda la generación de W24/S8 ni todos los
canales del estado terminal.

La ampliación coordinada \path{TerminalSelector}, de 18 de septiembre,
formaliza ya ese selector finito: 40.320 ordenaciones producen 19.446
candidatos distintos; 79 pertenecen a la cara excepcional y la condición
orbital deja un único superviviente. El módulo
\path{SelectedTerminalAlpha.lean} compone el registro seleccionado con una
raíz analítica única y con sus publicaciones dodecafásicas a cualquier
profundidad. Su recibo registra la compilación de cuatro módulos nuevos,
reutilizando las dependencias existentes. La enumeración utiliza
\path{native_decide} y, por ello, el informe declara
\path{Lean.ofReduceBool}; no se confunde esa comprobación con la reducción
exclusiva por el núcleo utilizada en los dos módulos locales de esta nota.

El módulo coordinado \path{SelectedFromRegionalInputs.lean} compone ahora
el selector con el descriptor, el panel y las firmas regionales generados.
El teorema \path{regional_terminal_alpha} conserva esa producción hasta la
raíz única y sus publicaciones a toda profundidad. El repertorio terminal
no ordenado mantiene una interfaz explícita, con su procedencia conservada
en el paquete; la nueva composición no presenta su producción anterior
como un teorema añadido. El recibo
\path{REGIONAL_COMPOSITION_VERIFICATION.json} registra los dos módulos de
composición y las huellas de sus dependencias regionales.

La composición coordinada posterior \path{SelectedRegionalIncidence.lean}
reúne ambas lecturas del mismo registro seleccionado. El soporte
$\{i:K_i\geq729\}$ coincide con la tétrada generada por los soportes
regionales. El soporte negativo del preacarreo
$P_i+E_i-\Phi_i-K_i$ coincide con la hexada seleccionada. Desde estas
igualdades, el módulo recupera el origen y el radial del retículo marcado,
con norma $54$, rango $24$, integralidad, paridad, autodualidad y mínimo
$4$ alcanzado. El teorema
\path{arithmetic_and_incidence_of_same_register} conserva conjuntamente
esta realización incidencial y la raíz aritmética con sus publicaciones
a toda profundidad. Las doce consultas de axiomas declaran la evaluación
nativa heredada del selector. El soporte de bloques altos y el signo del
preacarreo son los lectores efectivos; no se sustituyen por una reducción
módulo tres ni por el signo de una transformada distinta.

\section{Centro electrónico, bloque aritmético y registros regionales}

El artículo X ya reúne estas relaciones en
\path{sections/centro_electronico_registros.tex}. La celda $(5,5)$ es el
único punto fijo de $\rho(r,c)=(10-r,10-c)$. Sus hojas aritméticas dan
\[
 5+5=1+9\cdot1,\qquad 5\cdot5=7+9\cdot2.
\]
Sobre la misma fibra de suma diez, las posiciones conjugadas $(8,2)$ y
$(2,8)$ conservan el residuo multiplicativo siete y cambian el cociente
de dos a uno. La vacancia conserva así un defecto unitario identificable,
junto con la orientación de la posición alcanzada.

El bloque central evaluado sobre las posiciones adyacentes $4,5$ es
\[
 B_c=\begin{pmatrix}7&2\\2&7\end{pmatrix}=9P_++5P_-.
\]
Sus dos modos y sus lecturas diagonales pertenecen a esta misma
arquitectura; la posición $(5,5)$, el bloque matricial y los autovalores
son objetos relacionados mediante los lectores, no notaciones
intercambiables. La concatenación conserva
$72+27=77+22=99$ y $77-22=55=54+1$.

En la ruta electrónica, la orientación inicial $(1,1)$ y las ventanas de
nueve pasos determinan el registro
\[
 \tau_e=(+1,+1,+1,-1,-1,-1,+1,+1,+1,-1,-1,-1).
\]
La firma regional $555555$ procede, en cambio, de la emisión transversal
\[
 (501,614,498,169,272,272),
\]
mediante las diferencias decimales
$[d_2-d_1]_{10}$. La fuente demuestra la conexión con la palabra
regional $010211$ y conserva la memoria necesaria para transportar esa
firma: su fibra se divide en tres clases estables bajo los operadores
de avance y retorno. El texto y su figura se incluyen íntegros en el
paquete, junto con la recurrencia electrónica y sus lectores.

Esta articulación mantiene la jerarquía de la serie. La generación
regional y el registro de acoplamiento preceden a sus lecturas de
$\alpha$; los caracteres de acción, la escala determinantal y los
ángulos preceden a las aplicaciones que los utilizan. El funcional
hexada--octada, el área, la entropía, la gravitación y las realizaciones
electromagnéticas conservan sus desarrollos específicos en las fuentes
completas. La compilación focal de K no se confunde con una nueva
formalización Lean de todos esos sectores.

\section{Conservación evolutiva de la unidad y recuperación}

El artículo X contiene la cadena aritmética
\[
 (72,27)\longrightarrow(73,27)\longrightarrow(729,271).
\]
La primera operación incorpora la contribución de frontera registrada;
la segunda multiplica la escala por diez y transfiere una unidad entre
las dos coordenadas:
\[
 L_1(73,27)=(10\cdot73-1,10\cdot27+1)=(729,271).
\]
La pareja inicial de suma 99 no se declara por sí sola de suma cien:
la contribución de frontera pertenece a la información que la operación
conserva. Para $B=b+1$ y $p_b=b^2-b+1$, las identidades
\[
 Bp_b=b^3+1,\qquad p_b+3b=B^2
\]
producen
\[
 (p_b,3b)\longmapsto(Bp_b-1,3bB+1)=(b^3,B^3-b^3).
\]
La suma normalizada permanece igual a uno. El caso $b=9$ conserva
$271=243+27+1$ como descomposición del complemento.

Para un acarreo canónico, la transformación
\[
 L_c(x,y)=(Bx-c,By+c)
\]
tiene el recuperador
\[
 c=Y\bmod B,\qquad y=(Y-c)/B,\qquad x=(X+c)/B.
\]
Su iteración conserva la profundidad y el registro ordenado
$C_n=\sum_{j=1}^n c_jB^{n-j}$. El artículo demuestra la recuperación de
toda la palabra canónica, las condiciones adicionales para acarreos
firmados y la compatibilidad de los cilindros. El paquete conserva
íntegramente esas pruebas, no sólo las identidades anteriores.

La linealización sobre etiquetas admisibles produce el transporte
unitario J del artículo. Su composición bilateral conserva la norma y
transporta los lectores. En el caso nonádico, el lector comprimido cumple
\[
 W_8^*\operatorname{diag}(F,F)W_8
 =\frac{72F_J+R^*F_JR}{73},\qquad F_J=J^*FJ.
\]
La suma aritmética normalizada, la norma de amplitudes y la incidencia
excepcional son lecturas relacionadas por mapas explícitos, con sus
portadores respectivos. La recuperación de K conserva el marco necesario
para su lectura excepcional; la prolongación hacia Moonshine y las
realizaciones dimensionales permanece desarrollada en las secciones
completas del artículo que acompañan esta nota.

\section{Archivos e integración}

\begin{itemize}
\item \texttt{GeneratedMarkedIncidence.lean}: operadores reconstruidos,
calendario, semillas, incidencia marcada, origen y propiedades del retículo.
\item \texttt{WeightedIncidenceRecovery.lean}: testigos de doce hexadas,
inversa racional exacta, recuperación bilateral e inyectividad de la lectura
completa.
\item \texttt{verify\_lean.py}: clausura de importaciones, compilación y
registro de los axiomas utilizados por los teoremas consultados.
\item \texttt{control\_incidencia.py}: comprobaciones racionales exactas,
censo de hexadas y control del alcance de los soportes no ponderados.
\end{itemize}

Las fuentes textuales cotejadas son:
\begin{itemize}
\item Artículo X: secciones \path{generacion}, \path{registro_k},
\path{registro_incidencias}, \path{censo_terminal_completo},
\path{excepcional}, \path{fibras_app_witt} y \path{alpha}.
\item Integral: unidad U008 de órbitas, elevación y frontera R36.
\item Cierre afín--monodrómico terminal:
\path{TEOREMA_CIERRE_AFIN_MONODROMICO_TERMINAL.md}.
\item Procedencia del repertorio S8:
\path{DICTAMEN_PROCEDENCIA_SUPERVIVENCIA_S8.md}.
\end{itemize}
Las rutas y huellas se conservan en los recibos de la carpeta. Las
declaraciones del tratado completo no se consideran nuevamente certificadas
por la comprobación focal.

\end{document}
